\documentclass[12pt]{article}
\usepackage[utf8]{inputenc}
\usepackage[T1]{fontenc}
\usepackage[margin=1in]{geometry}
\usepackage{setspace}
\usepackage{xcolor}
\usepackage[colorlinks=true, urlcolor=blue!60!black, linkcolor=blue!60!black]{hyperref}
\setlength{\parindent}{2em}
\setlength{\parskip}{1.25em}
\renewcommand{\baselinestretch}{1.0}
\begin{document}
\begin{center}
{\LARGE \textit{Marsh v.\ Alabama}, 326 U.S.\ 501 (1946)\par}
\vspace{0.8em}
{\large Case Brief --- Four-Page Edition\par}
\end{center}
\vspace{1em}

\section{Memory Jogger}

A private corporation that owns and runs an entire company town open to the public performs a public function and is bound by the First and Fourteenth Amendments, so Alabama could not convict a Jehovah's Witness of trespass for leafletting on the town's business-block sidewalk (\href{https://supreme.justia.com/cases/federal/us/326/501/}{Justia}).

\section{Detailed Case Facts}

Chickasaw, a suburb of Mobile, was wholly owned by the Gulf Shipbuilding Corporation. Apart from private title, it had every attribute of a municipality: homes, streets, a sewer system, a sewage plant, and a ``business block'' of rented stores. A county deputy sheriff paid by Gulf served as police, and the United States ran a post office on the business block. Company streets connected to a four-lane public highway thirty feet away, and nothing kept travelers from entering. Grace Marsh, a Jehovah's Witness, stood on the company sidewalk near the post office distributing religious literature. Gulf had posted a notice that ``Without Written Permission, No Street, or House Vendor, Agent or Solicitation of Any Kind Will Be Permitted.'' Told she needed a permit that would not be issued, Marsh refused to leave.

\section{Procedural History}

The company-paid deputy arrested Marsh, and Alabama prosecuted her under a statute criminalizing remaining on another's premises after a warning. The trial court rejected her First and Fourteenth Amendment defense. The Alabama Court of Appeals affirmed because Gulf held title and had not dedicated the sidewalk to the public; the Alabama Supreme Court denied certiorari. The U.S.\ Supreme Court reversed on January 7, 1946.

\section{Judicial Votes}

\textbf{5--3}, reversed and remanded.
\begin{itemize}
\item \textbf{Majority:} Justice Black, joined by Justices Douglas, Murphy, and Rutledge; Justice Frankfurter joined except as to the Commerce Clause analogy.
\item \textbf{Concurring:} Justice Frankfurter.
\item \textbf{Dissenting:} Justice Reed, joined by Chief Justice Stone and Justice Burton.
\item \textbf{Not participating:} Justice Jackson.
\end{itemize}

\section{Holding}

A State may not, consistently with the freedoms of religion and the press, impose criminal punishment under a general trespass statute on a person for distributing religious literature on the sidewalk of a company town's freely accessible business district. A corporation that governs a community with all the attributes of a municipality performs a ``public function,'' and its property rights cannot override the ``preferred position'' of First Amendment liberties.

\section{Analysis of Opinions}

\textbf{Majority (Black).} \textit{Municipal baseline:} an ordinary town could not ban leafletting or demand a discretionary permit (\textit{Lovell v.\ Griffin}; \textit{Martin v.\ Struthers}). \textit{Title not dispositive:} state property law on ``dedication'' could not control the federal question; an owner who opens property to the public accepts limits. \textit{Public function:} like privately owned bridges, ferries, turnpikes, and railroads, facilities run for the public perform ``essentially a public function.'' \textit{Self-government:} millions lived in company towns, and as free citizens they needed uncensored information. Balancing property against freedom of press and religion, the latter prevail.

\textbf{Frankfurter (concurring).} Under the Court's Jehovah's Witness precedents, constitutional rights cannot turn on technical notions of trespass, because ``a company-owned town is a town.'' But the First Amendment is ``a denial of all governmental power,'' so the Commerce Clause analogy was inapt.

\textbf{Reed, joined by Stone and Burton (dissenting).} This ``novel Constitutional doctrine'' opens private property to unwanted speakers and requires ``new arbitrary lines, judicially drawn''; the utility analogy is inapt; the Constitution protects property too; and Marsh could have leafletted on the public highway thirty feet away.

\textbf{Later developments.} \textit{Logan Valley} (1968) extended \textit{Marsh} to shopping centers over Black's dissent, but \textit{Lloyd Corp.} (1972) and \textit{Hudgens} (1976) confined it to company towns. \textit{Jackson v.\ Metropolitan Edison} (1974) required functions ``traditionally exclusively reserved to the State,'' and \textit{Halleck} (2019) held that merely hosting speech is not such a function.

\section{Examples: Future Applications}

\begin{enumerate}
\item \textbf{Same-Side: Corporate Tech Town.} A company town with housing, utilities, open roads, a post office, and deputized security cannot have a labor leafletter arrested for trespass.
\item \textbf{Same-Side: Private Election Contractor.} A company running a county's elections is a state actor; elections are a traditionally exclusive public function (\textit{Terry v.\ Adams}).
\item \textbf{Opposite-Side: Social-Media Platform.} A dominant platform that suspends an activist is no state actor under \textit{Halleck}, and its curation is protected expression (\textit{Moody v.\ NetChoice}).
\item \textbf{Opposite-Side: Lifestyle Mall.} An open-air shopping center with faux ``Main Street'' sidewalks but no residences or services is not a town (\textit{Lloyd}; \textit{Hudgens}).
\item \textbf{Fence-Sitter: Developer-Owned New Town.} A private town whose police, fire, water, and schools are public looks like Chickasaw, but the Court has never said how many municipal attributes suffice.
\end{enumerate}

\section{Critique}

\textbf{Two theories blended.} \textit{Marsh} mixes the public-function theory with state enforcement through a deputy sheriff and a trespass statute, a precursor to \textit{Shelley v.\ Kraemer}; modern doctrine holds that neutral trespass enforcement is not state action, so the public-function label did the decisive work. Black also framed the result as balancing, while modern doctrine treats state action as a binary threshold.

\textbf{Obsolete or revived?} Industrial company towns have largely vanished, but master-planned communities, large HOAs, private campuses, and corporate ``innovation cities'' revive the question of how many municipal attributes turn private property into a de facto town.

\section{Key Quotations}

\begin{quote}
``Ownership does not always mean absolute dominion. The more an owner, for his advantage, opens up his property for use by the public in general, the more do his rights become circumscribed by the statutory and constitutional rights of those who use it.''
\end{quote}
\noindent --- Justice Black, for the Court, 326 U.S.\ at 506.

\begin{quote}
``Whether a corporation or a municipality owns or possesses the town, the public in either case has an identical interest in the functioning of the community in such manner that the channels of communication remain free.''
\end{quote}
\noindent --- Justice Black, for the Court, 326 U.S.\ at 507.

\begin{quote}
``When we balance the Constitutional rights of owners of property against those of the people to enjoy freedom of press and religion, as we must here, we remain mindful of the fact that the latter occupy a preferred position.''
\end{quote}
\noindent --- Justice Black, for the Court, 326 U.S.\ at 509.

\end{document}
